% answers_short.tex —— 2016 年·天津卷·文科 · 简短答案速查
% 由 tools/make_short_answers.py 自 answers.tex 的 \gksolhead/\gkanswer 提取，
% 并与 2016/2016_数学真题_天津文_答案要点.md 逐题比对后生成，勿手改（改源 MD 后重跑）。

\begin{gknotice}
\textbf{使用说明}：本节只给最终答案便于快速核对，完整推导见解析版。答案为编写者独立推导，\textbf{非官方标准答案}。
\end{gknotice}

\gkpart{选择题}
\begin{center}\begin{tabular}{ccccccccc}
\toprule
题号 & 1 & 2 & 3 & 4 & 5 & 6 & 7 & 8 \\
\midrule
答案 & A & A & B & A & C & C & B & D \\
\bottomrule
\end{tabular}\end{center}

\gkpart{填空题}
\begin{itemize}[leftmargin=2.6em]
  \item[\textbf{9.}] \(1\)
  \item[\textbf{10.}] \(3\)
  \item[\textbf{11.}] \(4\)
  \item[\textbf{12.}] \((x-2)^{2}+y^{2}=9\)
  \item[\textbf{13.}] \(\dfrac{2\sqrt{3}}{3}\)
  \item[\textbf{14.}] \(\left[\dfrac{1}{3},\dfrac{2}{3}\right)\)
\end{itemize}

\gkpart{解答题}
\begin{itemize}[leftmargin=2.6em]
  \item[\textbf{15.}] (1) \(B=\dfrac{\pi}{6}\)；(2) \(\sin C=\dfrac{1+2\sqrt{6}}{6}\)
  \item[\textbf{16.}] (1) \(4x+5y\leqslant 200\)，\(8x+5y\leqslant 360\)，\(3x+10y\leqslant 300\)，\(x\geqslant 0\)，\(y\geqslant 0\)；平面区域为该不等式组的解集（以 \((0,0)\)、\((45,0)\)、\((40,8)\)、\((20,24)\)、\((0,30)\) 为顶点的五边形区域）；(2) 甲 \(20\) 车皮、乙 \(24\) 车皮，最大利润 \(112\) 万元
  \item[\textbf{17.}] (1) 见解析版；(2) 见解析版；(3) \(\dfrac{\sqrt{5}}{6}\)
  \item[\textbf{18.}] (1) \(a_{n}=2^{n-1}\)；(2) \(2n^{2}\)
  \item[\textbf{19.}] (1) \(\dfrac{x^{2}}{4}+\dfrac{y^{2}}{3}=1\)；(2) 直线 \(l\) 的斜率 \(k=\pm\dfrac{\sqrt{6}}{4}\)
  \item[\textbf{20.}] (1) 当 \(a\leqslant 0\) 时单调递增区间为 \((-\infty,+\infty)\)、无单调递减区间；当 \(a>0\) 时单调递增区间为 \(\left(-\infty,-\sqrt{\dfrac{a}{3}}\right]\) 与 \(\left[\sqrt{\dfrac{a}{3}},+\infty\right)\)，单调递减区间为 \(\left[-\sqrt{\dfrac{a}{3}},\sqrt{\dfrac{a}{3}}\right]\)；(2) 见解析版；(3) 见解析版
\end{itemize}
